% =====================================================================
% Section 15, The three hypothesis tests (humanised; bundled
% notebook_sections figures)
% =====================================================================
\makesectioncover{The Three Hypothesis Tests}

\section{The Three Hypothesis Tests}
\label{sec:phase2-hypotheses}

This section is the formal evidence layer of the project. It opens with
a methods reference (\S\ref{sec:hyp-methods}) that explains every
statistical test in one place, and then walks through H1, H2 and H3
(\S\ref{sec:h1}, \S\ref{sec:h2}, \S\ref{sec:h3}). A reader who only
wants the narrative can skip the reference. A reader who wants to
defend any single number should read both.

% =====================================================================
\subsection{Methods Reference}
\label{sec:hyp-methods}

For every test we cover four things: what it is used for, why we
chose it over the parametric alternative, how it is calculated step
by step, and how to read the result.

% --------------------------------------------------------------------
\subsubsection{Spearman rank correlation \texorpdfstring{($\rho$)}{}}

\textbf{Used for.} Strength and direction of a monotonic association
between two continuous variables.

\textbf{Why over Pearson.} Pearson assumes a linear relationship and
is sensitive to outliers and to non-normal distributions. Spearman
operates on ranks, so it does not assume linearity, only monotonicity,
it is robust to outliers, and it does not need normality. Most of our
relationships (population vs. underservedness, density vs. informal
share) are monotone but not linear.

\textbf{How calculated.}

\begin{enumerate}
 \item Rank $x$ and $y$ separately, with average ranks for ties:
 $R(x_i)$ and $R(y_i)$.
 \item Compute differences $d_i = R(x_i), R(y_i)$.
 \item Apply the rank-correlation formula:
 \begin{equation*}
 \rho \;=\; 1 - \dfrac{6 \sum_{i=1}^{n} d_i^2}{n(n^2 - 1)}.
 \end{equation*}
 \item Significance via the t-approximation
 $t = \rho \sqrt{(n-2)/(1-\rho^2)}$ with $df = n-2$.
\end{enumerate}

\textbf{Worked example.} Take $x = (1, 2, 3, 4, 5)$ and
$y = (2, 1, 4, 3, 5)$. Then $R(x) = (1, 2, 3, 4, 5)$, $R(y) = (2, 1, 4, 3, 5)$,
$d = (-1, 1, -1, 1, 0)$, $\sum d^2 = 4$.
$\rho = 1, \tfrac{6 \cdot 4}{5(24)} = 1, 0.2 = 0.8$.

\textbf{Reading.} $|\rho| < 0.1$ negligible, $0.1$ to $0.3$ weak,
$0.3$ to $0.5$ moderate, more than $0.5$ strong. Q17's
$\rho = 0.548$ at $p < 10^{-100}$ is strong; Q18's $\rho = 0.025$ at
$p = 0.40$ is statistically zero.

% --------------------------------------------------------------------
\subsubsection{Kruskal-Wallis H test}

\textbf{Used for.} Comparing the central tendency of more than two
independent groups when normality cannot be assumed.

\textbf{Why over one-way ANOVA.} ANOVA assumes normal residuals
within each group and equal variance across groups. Stations per 100k
residents across density tertiles are heavy-tailed and have different
variances, so ANOVA's p-value would not be reliable. Kruskal-Wallis
ranks every observation globally, so distribution shape stops
mattering.

\textbf{How calculated.}

\begin{enumerate}
 \item Pool all $N = \sum_g n_g$ observations across $k$ groups and
 rank globally; $R_{ij}$ is the rank of observation $j$ in
 group $i$.
 \item Compute the per-group rank sum $R_g = \sum_j R_{gj}$.
 \item The test statistic is
 \begin{equation*}
 H \;=\; \frac{12}{N(N+1)} \sum_{g=1}^{k} \frac{R_g^2}{n_g} \;-\; 3(N+1).
 \end{equation*}
 \item Under $H_0$ (all groups share the same distribution),
 $H \sim \chi^2_{k-1}$ asymptotically. Reject $H_0$ when
 $H > \chi^2_{1-\alpha,\, k-1}$.
\end{enumerate}

\textbf{Reading.} A significant $H$ tells you that at least one group
differs. It does not say which, and it does not measure how much. To
say which, run pairwise Mann-Whitney with Holm correction. To say
how much, compute $\varepsilon^2$.

% --------------------------------------------------------------------
\subsubsection{Epsilon-squared \texorpdfstring{($\varepsilon^2$)}{} effect size}

\textbf{Used for.} The fraction of rank-variance explained by group
membership.

\textbf{Why this and not eta-squared.} $\eta^2$ assumes the original
sums of squares; we already moved to ranks, so the rank-equivalent
$\varepsilon^2$ is the apples-to-apples effect size.

\textbf{How calculated.}
\begin{equation*}
 \varepsilon^2 \;=\; \frac{H - k + 1}{N - k}.
\end{equation*}

(That is the convention scipy and pingouin report.)

\textbf{Reading.} Less than $0.01$ negligible, $0.01$ to $0.04$ small,
$0.04$ to $0.16$ medium, more than $0.16$ large. H1's
$\varepsilon^2 = 0.826$ sits roughly five times above the large
threshold, far outside any range we would expect from chance.

% --------------------------------------------------------------------
\subsubsection{Mann-Whitney U test}

\textbf{Used for.} Comparing two independent groups when the
distributions are non-normal, ordinal, or heavily skewed.

\textbf{Why over Welch's t-test.} Welch relaxes the equal-variance
assumption of the standard t-test, but keeps the normality assumption.
Mann-Whitney drops both and works on ranks. With small $n$ (H2 has
11 plus 31 stations, H3 has 12 plus 12), this matters a lot.

\textbf{How calculated.}

\begin{enumerate}
 \item Pool the two groups, rank globally, average ranks for ties.
 \item Compute $R_1$, the rank-sum of group 1.
 \item Derive $U_1 = R_1 - \tfrac{n_1(n_1+1)}{2}$,
 $U_2 = n_1 n_2 - U_1$, take $U = \min(U_1, U_2)$.
 \item Small-sample p-values use the exact Mann-Whitney distribution
 (scipy lookup). Large-sample p-values use the normal
 approximation
 \begin{equation*}
 z \;=\; \frac{U - \tfrac{n_1 n_2}{2}}
 {\sqrt{\tfrac{n_1 n_2 (n_1+n_2+1)}{12}}}.
 \end{equation*}
 \item Tied ranks subtract a known correction term from the variance.
\end{enumerate}

\textbf{Reading.} The p-value answers whether the two distributions
are stochastically different. It does not say how different, so we
always pair it with Cliff's delta.

% --------------------------------------------------------------------
\subsubsection{Cliff's delta \texorpdfstring{($\delta$)}{} effect size}

\textbf{Used for.} A non-parametric effect size that goes with a
Mann-Whitney result. It is the probability that a randomly drawn $x$
exceeds a randomly drawn $y$, minus the reverse.

\textbf{Why this and not Cohen's d.} Cohen's d uses a standardised
mean difference, which assumes normality and breaks under skew.
Cliff's $\delta$ counts pairwise dominance and is robust to outliers
and shape differences.

\textbf{How calculated.}
\begin{equation*}
 \delta \;=\; \frac{\#(x_i > y_j) \;-\; \#(x_i < y_j)}{n_1 \cdot n_2}.
\end{equation*}

The range is $[-1, +1]$. $\delta = +1$ means every $x$ exceeds every
$y$; $\delta = -1$ is the reverse; $\delta = 0$ is 50/50.

\textbf{Reading.} $|\delta| < 0.147$ negligible, $0.147$ to $0.33$
small, $0.33$ to $0.474$ medium, more than $0.474$ large. H2's
$\delta = -0.991$ means 99.55\% of LRT-vs.-metro pairwise comparisons
go the same way (metro larger). H3's $\delta = +0.710$ means BRT
corridors dominate matched controls in 85.5\% of pairwise
comparisons.

% --------------------------------------------------------------------
\subsubsection{Permutation test}

\textbf{Used for.} Building a null distribution empirically when the
test statistic has no clean closed-form distribution. We use it for
the Moran's I p-value and as a robustness check for H3.

\textbf{Why this and not bootstrap.} A permutation test preserves
marginals and tests exchangeability under $H_0$ directly. That is the
right tool for the question \emph{did this difference happen by
chance}. Bootstrap estimates the sampling variability of an
estimate, which is a different question.

\textbf{How calculated.}
\begin{enumerate}
 \item Compute the observed statistic $T_\text{obs}$ on the actual
 data.
 \item For $b = 1, \ldots, B$ iterations, shuffle the labels and
 recompute $T_b$.
 \item Empirical p-value:
 \begin{equation*}
 p \;=\; \frac{\#\{|T_b| \geq |T_\text{obs}|\} + 1}{B + 1}.
 \end{equation*}
\end{enumerate}

We use $B = 999$ for Moran's I and $B = 10{,}000$ for the H3
robustness check.

% --------------------------------------------------------------------
\subsubsection{Moran's I (spatial autocorrelation, AI Technique 3)}

\textbf{Used for.} Detecting whether values at nearby spatial units
resemble each other more than chance. In our case it answers whether
the H1 mismatch is geographically clustered or randomly scattered.

\textbf{How calculated.}

\begin{enumerate}
 \item Build a spatial weight matrix $W$. We use KNN-based weights
 (each district's 5 nearest neighbours by Nominatim centroid
 distance), row-standardised. Polygon adjacency is unavailable
 for these districts, so KNN acts as a contiguity proxy.
 \item Centre the variable: $z_i = x_i, \bar{x}$.
 \item Compute Moran's I:
 \begin{equation*}
 I \;=\; \frac{n}{\sum_{i,j} w_{ij}}
 \cdot \frac{\sum_{i,j} w_{ij} z_i z_j}
 {\sum_i z_i^2}.
 \end{equation*}
 \item Under $H_0$ (no spatial autocorrelation),
 $E[I] = -1/(n-1)$, slightly negative for finite $n$.
 \item Empirical p-value via 999 permutations of the $z$ vector
 against the fixed $W$.
\end{enumerate}

\textbf{Reading.} $I \approx E[I]$ means no clustering. $I > E[I]$
means similar values cluster (positive autocorrelation). $I < E[I]$
means dispersion. The H1 result for stations per 100k is positive
and significant once we use real Nominatim centroids.

% --------------------------------------------------------------------
\subsubsection{K-Means clustering (AI Technique 2)}

\textbf{Used for.} Partitioning $n$ observations into $k$ groups by
minimising the within-cluster sum of squared distances.

\textbf{How calculated.}
\begin{enumerate}
 \item Pick $k$.
 \item Initialise $k$ centroids using k-means++ (probabilistic seeding
 biased away from already-chosen centroids, which gives stable
 outputs from random seeds).
 \item Iterate: assign each point to its nearest centroid, recompute
 centroid as the mean of assigned points, until assignments
 stop changing.
 \item Validate by re-running with multiple seeds and computing the
 Adjusted Rand Index (ARI). High ARI means the clustering is
 reproducible.
\end{enumerate}

\textbf{Q24 settings.} $k = 4$, justified by elbow plus silhouette
of $0.412$, \fn{n\_init=50}, ARI of \stat{1.00} across five seeds.

% =====================================================================
\subsection{H1 \textbullet\ Coverage-need mismatch}
\label{sec:h1}

\subsubsection{Setup}

\textbf{Story-level claim.} Cairo's transport infrastructure penalises
density. Districts with more people per square kilometre receive fewer
stations per 100k residents than less-dense districts, which is the
opposite of what a demand-aligned planner would do.

\textbf{Hypotheses.}
\begin{align*}
 H_0:\;& \text{coverage per 100k residents is the same across density tertiles} \\
 H_1:\;& \text{at least one tertile differs}
\end{align*}

A secondary geographic hypothesis runs alongside the main test:
\begin{align*}
 H_0^{(\text{spatial})}:\;& \text{under-service is randomly scattered across the city} \\
 H_1^{(\text{spatial})}:\;& \text{under-service clusters geographically}
\end{align*}

\subsubsection{Methods chosen and why}

The main test needs to compare more than two groups and tolerate
non-normal distributions. That is exactly what Kruskal-Wallis is for (see
\S\ref{sec:hyp-methods}). $\varepsilon^2$ sizes the effect; pairwise
Mann-Whitney with Holm correction identifies which tertiles differ;
Moran's I answers the spatial-clustering question.

\subsubsection{Data pipeline trace}

\begin{enumerate}
 \item 68 Greater Cairo districts from S6, with real Nominatim
 centroids. (Earlier drafts used three governorate centroids
 and gave $\varepsilon^2 \approx 0.21$; the upgrade is what
 produced the large effect we report.)
 \item Population from S6's 2017 column.
 \item Station count is the number of integrated post-2014 stations
 within a 3 km buffer of each district centroid.
 \item Per-district metric: \fn{stations\_per\_100k = stations /
 (population / 100{,}000)}.
 \item Split into low, medium and high tertiles
 ($n_\text{low} = 23$, $n_\text{med} = 22$, $n_\text{high} = 23$).
 \item Run \fn{scipy.stats.kruskal(low, med, high)}.
 \item Run pairwise \fn{scipy.stats.mannwhitneyu} with Holm
 correction.
 \item Run \fn{esda.Moran(values, KNN(centroids, k=5), permutations=999)}.
\end{enumerate}

\subsubsection{Results}

Tertile medians (stations per 100k): low \stat{25.9}, medium
\stat{8.5}, high \stat{2.15}, a roughly 12$\times$ ratio. Kruskal-Wallis
\stat{$H = 55.7$, $p < 0.001$}, $\varepsilon^2 = \stat{0.826}$
(large). Pairwise Mann-Whitney with Holm correction: every pair
differs at $p < 0.001$. Moran's I on real Nominatim centroids
indicates positive spatial clustering of under-service.

\figitem{phase2/sections/h1_notebook_visuals.png}{1.0}{H1 bundled
visuals from the analysis notebook.}{fig:h1bundle}

\subsubsection{What the result means}

\begin{callout}
The 12$\times$ tertile ratio combined with $\varepsilon^2 = 0.826$ is
not a small effect that crossed a significance line. It is the
quantitative signature of a planning logic that built infrastructure
where future residents are expected to arrive (the New Administrative
Capital, the satellite cities) rather than where dense Cairo already
sits. The Moran's I result says under-service concentrates rather
than scatters.

The methodology upgrade (real Nominatim centroids) is mentioned to
be honest about the use. The same data with worse spatial
resolution would have told a noticeably weaker story.
\end{callout}

% =====================================================================
\subsection{H2 \textbullet\ LRT catchment deficit}
\label{sec:h2}

\subsubsection{Setup}

\textbf{Story-level claim.} The LRT was built where almost no one
currently lives. Each operational LRT station has a much smaller 2 km
catchment population than each post-2012 metro Line 3 station.

\textbf{Hypotheses.}
\begin{align*}
 H_0:\;& \text{LRT and post-2012 metro L3 have the same 2 km catchment population distribution} \\
 H_1:\;& \text{the two distributions differ}
\end{align*}

\subsubsection{Methods chosen and why}

Mann-Whitney U for the two-group comparison (non-normal catchments
with small $n$). Cliff's $\delta$ for the effect size, because a
p-value alone says nothing about how big the gap is.

\subsubsection{Data pipeline trace}

\begin{enumerate}
 \item 11 operational LRT stations from S4 (after the 9 Overpass
 plus 7 Google Maps coordinate backfills).
 \item 31 metro L3 post-2012 stations from S3 (filtered by opening
 year $> 2012$).
 \item Draw a 2 km buffer per station in EPSG:32636.
 \item Sum \fn{Population\_2018} over every Phase 1 H3 hex the
 buffer intersects.
 \item Run \fn{scipy.stats.mannwhitneyu(lrt\_pop, metro\_pop,
 alternative='two-sided')}.
 \item Compute Cliff's $\delta$ from the pairwise comparison matrix.
 \item Sensitivity: rerun at 1 km and 3 km radii and confirm
 $\delta$ stays below $-0.95$ at every radius.
\end{enumerate}

\subsubsection{Results}

Medians inside the 2 km buffer: LRT = \stat{916} residents, metro L3
post-2012 = \stat{634{,}333} residents. The ratio is roughly
\stat{688$\times$}. Mann-Whitney $U = \stat{2}$, $p < \stat{0.0001}$.
Cliff's $\delta = \stat{-0.991}$, within a hair of the $-1$ floor,
which means 99.55\% of pairwise comparisons go the same way.
Sensitivity holds at 1 km and 3 km radii.

\figitem{phase2/sections/h2_notebook_visuals.png}{1.0}{H2 bundled
visuals.}{fig:h2bundle}

\subsubsection{What the result means}

\begin{callout}
Cliff's $\delta = -0.991$ is the floor. The reading is
mechanical: the LRT was sited for a city that does not yet exist. The
planning logic was forward-looking: build the infrastructure first,
let settlement follow. That logic may pay off as the New Administrative
Capital fills, but for 2024-2026 the gap between the LRT and the
people who would ride it is the largest single mismatch in Phase 2.
\end{callout}

% =====================================================================
\subsection{H3 \textbullet\ BRT corridor match}
\label{sec:h3}

\subsubsection{Setup}

\textbf{Story-level claim.} The BRT, unlike the LRT and the post-2012
metro extensions, was sited on top of existing informal demand. The
Ring-Road BRT runs along corridors that already carry substantial
microbus boarding.

\textbf{Hypotheses.}
\begin{align*}
 H_0:\;& \text{BRT 500 m buffers and matched random non-Ring-Road controls have} \\
 & \text{the same informal demand distribution} \\
 H_1:\;& \text{BRT buffers carry more informal demand than controls}
\end{align*}

\subsubsection{Methods chosen and why}

Mann-Whitney U again, for the same reasons as H2. Cliff's $\delta$
for the effect size. The H3-specific contribution is the matched
control sample: random points drawn from the urbanised non-Ring-Road
area at similar distance to the city centre as the BRT corridor.

\subsubsection{Data pipeline trace}

\begin{enumerate}
 \item 12 BRT stations from S5.
 \item Per BRT station, draw a 500 m buffer and sum Phase 1 informal
 boardings for stops inside the buffer.
 \item Generate 12 random control points from the urbanised
 non-Ring-Road area and replicate the buffer-and-sum step.
 \item Run \fn{scipy.stats.mannwhitneyu(brt\_demand, ctrl\_demand,
 alternative='greater')}.
 \item Compute Cliff's $\delta$.
 \item Cross-validate with a 10{,}000-iteration permutation test.
\end{enumerate}

\subsubsection{Results}

Medians inside the 500 m buffer: BRT \stat{566} informal boardings,
control \stat{0}. Mann-Whitney $U = \stat{132}$, $p = \stat{0.0001}$.
Cliff's $\delta = \stat{+0.710}$, large positive. Permutation
cross-validation confirms the same conclusion.

\figitem{phase2/sections/h3_notebook_visuals.png}{1.0}{H3 bundled
visuals.}{fig:h3bundle}

\subsubsection{What the result means}

\begin{callout}
H3 is the contrast case. Same state planner, same \$10B programme,
opposite outcome from H2. The BRT corridor demand is structurally
higher than what a random non-Ring-Road sample would carry, which
means the BRT was routed through existing demand rather than next to
it. H3's job in this report is to keep us from over-claiming.
\emph{All new infrastructure is misaligned} is wrong; \emph{most new
infrastructure is misaligned, but the BRT is the demand-aligned
exception} is what the data actually says. The Q18b matrix has its
single 25\% cell at BRT $\times$ G3 for the same reason.
\end{callout}

% =====================================================================
\subsection{Summary table}
\label{sec:hyp-summary}

\rowcolors{2}{darkpanel}{darkbg}
\begin{table}[H]
\centering\small
\renewcommand{\arraystretch}{1.4}
\begin{tabular}{@{} >{\color{plotorange}\bfseries}l
 >{\color{bodytext}}l
 >{\color{plotyellow}\ttfamily}l
 >{\color{plotyellow}\ttfamily}l
 >{\color{plotyellow}\ttfamily}l @{}}
\toprule
\color{plotblue}H & \color{plotblue}Test & \color{plotblue}Statistic & \color{plotblue}p-value & \color{plotblue}Effect size \\
\midrule
H1 & Kruskal-Wallis (3 groups, $n=68$) & $H = 55.7$ & $< 0.001$ & $\varepsilon^2 = 0.826$ \\
 & Moran's I (999 perms) & positive & significant & --- \\
H2 & Mann-Whitney U ($n=11$ vs. 31) & $U = 2$ & $< 0.0001$ & $\delta = -0.991$ \\
H3 & Mann-Whitney U ($n=12$ vs. 12) & $U = 132$ & $0.0001$ & $\delta = +0.710$ \\
\bottomrule
\end{tabular}
\caption{Three hypotheses, one summary. H1 establishes the systemic
density penalty; H2 establishes the most extreme single-mode failure
(LRT); H3 establishes that the contrast exists (BRT). Together they
support the verdict that not every \$10B station is misaligned, but
most are.}
\end{table}
\rowcolors{0}{}{}
